\documentclass[12pt,a4paper]{article}

\usepackage[T1]{fontenc}
\usepackage[utf8]{inputenc}
\usepackage[brazil]{babel}
\usepackage{lmodern}
\usepackage{microtype}
\usepackage{geometry}
\usepackage{indentfirst}
\usepackage{setspace}
\usepackage{enumitem}
\usepackage{array}
\usepackage{booktabs}
\usepackage{tabularx}
\usepackage{longtable}
\usepackage{float}
\usepackage[hidelinks]{hyperref}
\usepackage{xurl}

\geometry{top=2cm,bottom=2cm,left=2cm,right=2cm}
\setstretch{1.15}
\setlength{\parindent}{1.25cm}
\setlength{\parskip}{0.10cm}
\renewcommand{\arraystretch}{1.15}
\setlist[itemize]{leftmargin=1.05cm,itemsep=0.08cm,topsep=0.10cm}

\newcolumntype{Y}{>{\raggedright\arraybackslash}X}
\newcolumntype{L}[1]{>{\raggedright\arraybackslash}p{#1}}

\hypersetup{
  pdftitle={Desenvolvimento de um Aplicativo Mobile para Conferência de Bens Patrimoniais Integrado ao EGAP},
  pdfauthor={Lucas César},
  pdfsubject={Gestão patrimonial pública e desenvolvimento mobile},
  pdfkeywords={EGAP, inventário patrimonial, aplicativo mobile, Laravel, React Native, API REST}
}

\title{\textbf{Desenvolvimento de um Aplicativo Mobile para Conferência de Bens Patrimoniais Integrado ao EGAP}}
\author{Lucas César\\Universidade Vila Velha}
\date{2025}

\begin{document}

\maketitle

\begin{abstract}
Este trabalho apresenta o desenvolvimento de uma solução mobile integrada ao
EGAP, sistema de gestão administrativa e patrimonial utilizado em contexto do
Poder Público, com foco na conferência de bens móveis durante inventários. O
problema abordado decorre da separação entre a vistoria física, realizada nos
setores, e o registro digital, tradicionalmente dependente de anotações ou
acesso posterior a computador. A solução proposta é composta por um aplicativo
desenvolvido com Expo e React Native e por uma API construída em Laravel,
autenticada com Laravel Sanctum e integrada à base relacional legada. O estudo
caracteriza-se como pesquisa aplicada, exploratória e descritiva, conduzida por
meio de estudo de caso tecnológico e análise da implementação. Foram
implementadas funcionalidades de autenticação, consulta patrimonial, leitura
por câmera ou digitação, listagem dos bens vinculados ao setor, registro de
localização, não localização e divergência, acompanhamento de métricas e
finalização controlada da atividade. As regras sensíveis permanecem no
backend, que resolve o escopo do servidor autenticado, valida estados e grava
alterações em transação. Como resultado preliminar, a aplicação fornece uma
base funcional para aproximar o registro da conferência do local em que o bem
é verificado, aumentar a rastreabilidade do processo e reduzir etapas manuais,
sem substituir os controles institucionais existentes.

\noindent\textbf{Palavras-chave:} gestão patrimonial; inventário; aplicativo
mobile; Poder Público; Laravel; React Native; API REST.
\end{abstract}

\tableofcontents
\clearpage

\section{Introdução}

A gestão de bens permanentes no setor público exige que informações sobre
identificação, localização e responsabilidade estejam disponíveis para
controle, prestação de contas e tomada de decisão. A Lei n.º 4.320/1964
determina a existência de registros analíticos dos bens de caráter permanente
e estabelece que o levantamento dos bens móveis e imóveis tenha por base o
inventário analítico de cada unidade administrativa \cite{brasil4320}. A NBC
TSP 07, por sua vez, disciplina o ativo imobilizado no setor público e reforça
a importância de identificação, mensuração e controle dos ativos
\cite{cfcnbctsp07}.

Nesse domínio, o EGAP (Sistema de Gestão Administrativa Patrimonial) é uma
aplicação institucional que reúne operações de cadastro, patrimônio, pedidos,
almoxarifado, relatórios e administração. O sistema existente utiliza Laravel
11 e Filament 3 na camada web administrativa e mantém dados patrimoniais em
banco relacional MySQL/MariaDB. Entre suas estruturas estão os bens móveis, os
ciclos de inventário, as atividades por setor e os itens efetivamente
registrados durante a conferência.

Embora o sistema desktop forneça o núcleo administrativo, a conferência física
de bens é uma atividade de campo. O servidor precisa circular pelo setor,
localizar etiquetas, comparar características observadas e registrar
ocorrências. Quando a inserção dos dados depende de papel, planilha ou retorno
a uma estação de trabalho, cria-se um intervalo entre a observação física e o
registro digital. Esse intervalo favorece retrabalho, esquecimento de
justificativas, duplicidade de registros e menor visibilidade sobre o progresso
do inventário.

Este artigo descreve o desenvolvimento do \textit{Inventário Mobile}, módulo
mobile integrado ao EGAP para apoiar a conferência setorial de bens
patrimoniais. O aplicativo foi construído em Expo/React Native e se comunica
com uma API Laravel por HTTP, utilizando autenticação por token com Laravel
Sanctum. A implementação prioriza uma decisão arquitetural: o celular é a
interface operacional de campo, mas o backend continua responsável por
autenticação, escopo, validações patrimoniais, transações e persistência.

\subsection{Descrição e formulação do problema}

O problema investigado consiste na ausência de uma interface móvel integrada ao
EGAP para que servidores executem a conferência patrimonial no próprio ambiente
em que os bens se encontram. Sem esse suporte, a verificação física e a
atualização do inventário podem ocorrer em momentos ou meios distintos,
dificultando a rastreabilidade das ações e a atualização imediata dos
indicadores do setor.

Assim, formula-se a seguinte questão de pesquisa: \textit{como desenvolver uma
solução mobile integrada ao EGAP que permita realizar a conferência de bens
patrimoniais em campo, preservando o escopo do usuário autenticado, as regras
do sistema legado e a integridade dos registros de inventário?}

\subsection{Objetivos}

O objetivo geral deste trabalho é desenvolver uma solução mobile integrada ao
EGAP para apoiar a conferência de bens móveis em inventários patrimoniais
setoriais.

Como objetivos específicos, definem-se:

\begin{itemize}
  \item analisar o processo de conferência e as estruturas patrimoniais já
  existentes no EGAP;
  \item definir uma arquitetura cliente-servidor que mantenha as regras
  críticas no backend institucional;
  \item implementar autenticação mobile e resolução segura do setor e da
  unidade do usuário;
  \item implementar consulta, leitura e registro de estados de conferência por
  meio de uma API;
  \item disponibilizar uma interface mobile para consulta, indicadores,
  conferência e encerramento da atividade;
  \item analisar qualitativamente os resultados iniciais, limitações e
  possibilidades de evolução da solução.
\end{itemize}

\subsection{Justificativa e delimitação}

A proposta é relevante por tratar uma necessidade operacional diretamente
relacionada ao controle de patrimônio público: registrar com maior proximidade
temporal e espacial aquilo que foi observado na vistoria. A mobilidade permite
que a leitura de um patrimônio, o apontamento de divergência e a consulta de
pendências ocorram durante a conferência, enquanto a integração com o sistema
existente preserva regras e dados já utilizados institucionalmente.

O trabalho concentra-se na conferência de bens móveis vinculados à unidade e ao
setor do usuário autenticado. Fazem parte do recorte: autenticação, consulta,
listagem, leitura de código, registro de localização, não localização,
divergência, métricas e finalização da atividade. A primeira versão não tem
como objetivo executar transferência patrimonial definitiva, alterar
automaticamente dados cadastrais divergentes, operar em modo offline,
disponibilizar assinatura eletrônica ou avaliar quantitativamente ganhos de
tempo em produção.

\section{Referencial Teórico}

\subsection{Patrimônio público e inventário de bens móveis}

No setor público, o patrimônio permanente precisa ser controlado tanto sob a
perspectiva contábil quanto sob a perspectiva administrativa. Os artigos 94 e
96 da Lei n.º 4.320/1964 associam a caracterização individual do bem, a
responsabilidade por sua guarda e o inventário da unidade administrativa
\cite{brasil4320}. A conferência física, portanto, não é mera listagem: ela
apoia a confirmação da existência, da localização e das condições dos ativos
registrados.

A NBC TSP 07 define o ativo imobilizado como ativo tangível mantido para uso na
produção ou fornecimento de bens e serviços, para fins administrativos ou em
outras hipóteses previstas pela norma. Seu tratamento envolve reconhecimento,
mensuração e depreciação, o que depende de registros confiáveis ao longo da
vida do bem \cite{cfcnbctsp07}. Uma aplicação para conferência não substitui o
tratamento contábil, mas melhora a coleta da informação operacional que o
suporta.

Nesse contexto, estados como localizado, não localizado, divergente ou em
transferência são úteis para representar ocorrências verificadas durante o
inventário. O registro de justificativa e de autoria das alterações permite que
uma ocorrência seja posteriormente analisada por responsáveis administrativos,
sem confundir a observação em campo com uma alteração definitiva do cadastro
mestre.

\subsection{Arquitetura de software, MVC e camada de serviço}

Arquitetura de software corresponde à organização de elementos de uma solução,
às responsabilidades atribuídas a esses elementos e às relações necessárias
para atender aos atributos de qualidade do sistema. Em uma aplicação
institucional, manutenibilidade, confiabilidade, segurança e evolução são
especialmente relevantes, pois o software precisa conviver com regras
administrativas e dados persistentes \cite{bass2021}.

O padrão \textit{Model-View-Controller} (MVC) separa a representação dos
dados, a interação do usuário e a coordenação das requisições. No EGAP, os
\textit{models} Eloquent representam entidades persistidas; os
\textit{controllers} da API recebem requisições e formatam respostas; e a
interface administrativa ou o aplicativo mobile apresentam informações e
ações. O Laravel fornece roteamento, injeção de dependências, validação e o
Eloquent ORM para essa organização \cite{laravelrouting,laraveleloquent}.

Para a conferência, apenas a separação controller/model seria insuficiente,
pois localizar inventário, calcular pendências, impedir duplicidade e efetuar
gravações consistentes formam regras de negócio reutilizáveis. O padrão
\textit{Service Layer} organiza essas operações em um serviço de aplicação que
define a fronteira do caso de uso e coordena respostas e persistência
\cite{fowlerServiceLayer}. A implementação analisada concretiza esse princípio
no \texttt{ConferenciaBensService}, evitando que a regra patrimonial fique
dispersa nas telas mobile.

\subsection{Banco de dados relacional e transações}

O modelo relacional, introduzido por Codd, representa informações por relações
e busca independência entre a organização lógica dos dados e os detalhes de
armazenamento \cite{codd1970}. Tal abordagem é apropriada ao inventário
patrimonial porque bens, inventários, atividades, unidades e registros de
conferência possuem vínculos estruturados e precisam ser consultados de forma
consistente.

O MySQL é um sistema gerenciador de banco de dados relacional; com o mecanismo
InnoDB, disponibiliza transações e mecanismos de integridade aplicáveis a
operações que não devem permanecer parcialmente gravadas
\cite{mysqlTransactions,mysqlForeignKeys}. No EGAP, confirmar a localização de
um bem envolve criar um item de inventário, atualizar campos legados do bem e
recalcular a quantidade inventariada. Essas alterações são executadas em
transação na conexão institucional. Conforme a documentação do Laravel, a
operação \texttt{DB::transaction} efetiva o conjunto de operações apenas em
caso de sucesso e realiza \textit{rollback} em caso de exceção
\cite{laraveldatabase}.

\subsection{API REST, HTTP e HTTPS}

REST (\textit{Representational State Transfer}) é um estilo arquitetural para
sistemas distribuídos descrito por Fielding, no qual recursos são manipulados
por representações e por uma interface uniforme, favorecendo separação entre
cliente e servidor \cite{fielding2000}. O aplicativo do presente trabalho não
acessa diretamente as tabelas patrimoniais; ele solicita recursos e ações à
API, como consultar bens ou registrar uma ocorrência de conferência.

HTTP é um protocolo de aplicação sem estado baseado na troca de requisições e
respostas. A RFC 9110 define os métodos, códigos de resposta e semântica
compartilhada pelas versões modernas do protocolo \cite{rfc9110}. Na API do
EGAP, métodos \texttt{GET} são utilizados em consultas e métodos
\texttt{POST} registram operações de domínio como localizar, declarar
divergência ou finalizar atividade. Erros de validação, autorização ou conflito
podem ser comunicados ao cliente por respostas HTTP correspondentes.

Em implantação real, a comunicação que inclui credenciais e tokens deve ocorrer
por HTTPS. O esquema HTTPS identifica comunicação HTTP protegida por TLS; o TLS
1.3 foi padronizado para proteger a comunicação contra interceptação e
adulteração em trânsito \cite{rfc9110,rfc8446}. O uso de HTTPS complementa,
mas não substitui, autenticação, autorização e controle de escopo no backend.

\subsection{Autenticação por token e escopo de acesso}

Laravel Sanctum oferece autenticação para aplicações mobile e APIs simples por
meio de tokens enviados no cabeçalho \texttt{Authorization}. Tokens podem ser
emitidos e revogados, e rotas podem ser protegidas pelo middleware
\texttt{auth:sanctum} \cite{laravelsanctum}. Essa abordagem permite que a API
permaneça sem sessão de navegador e reconheça o usuário que realizou cada
requisição.

No domínio patrimonial, autenticar não é suficiente. O sistema precisa
determinar em qual unidade e setor o servidor pode conferir bens. Assim, o
EGAP resolve o escopo no backend a partir do usuário vinculado e de sua
lotação; o aplicativo não envia setor ou unidade como autoridade para a
operação. Essa decisão reduz a possibilidade de apontamentos feitos
arbitrariamente sobre bens de outro setor.

\subsection{Aplicações mobile com React Native e Expo}

React Native permite compor interfaces para plataformas móveis utilizando
componentes React que interagem com elementos nativos das plataformas
\cite{reactnativeDocs}. No projeto, TypeScript adiciona tipagem aos contratos
de API e aos estados da interface, contribuindo para explicitar os formatos
esperados nas operações de conferência.

Expo fornece bibliotecas e ferramentas sobre React Native. O Expo Router
organiza a navegação por arquivos: telas armazenadas no diretório
\texttt{app} tornam-se rotas do aplicativo, modelo utilizado para login,
painel, listagem e conferência \cite{expoRouter}. O pacote
\texttt{expo-camera} viabiliza leitura de códigos por câmera
\cite{expoCamera}; já \texttt{expo-secure-store} armazena pares chave-valor de
forma protegida no dispositivo, utilizando mecanismos da plataforma móvel
\cite{expoSecureStore}. No aplicativo desenvolvido, o token de acesso é
armazenado por essa biblioteca em Android e iOS; na variante web de
desenvolvimento, há uma estratégia distinta baseada em armazenamento do
navegador, que não deve ser tratada como equivalente de segurança para
produção mobile.

\section{Metodologia}

\subsection{Natureza e abordagem da pesquisa}

Este trabalho caracteriza-se como pesquisa aplicada, pois busca solucionar um
problema operacional concreto: apoiar digitalmente a conferência de bens
patrimoniais em campo. Possui caráter exploratório e descritivo, uma vez que
analisa um sistema existente, identifica seus fluxos e descreve a construção de
um módulo integrado. A abordagem é qualitativa, pois os resultados desta etapa
são analisados pela cobertura funcional e pelas decisões arquiteturais
implementadas, sem experimento estatístico com usuários finais.

Como estratégia, adotou-se um estudo de caso tecnológico sobre o EGAP e o
Inventário Mobile. O objeto do estudo é a conferência setorial de bens móveis,
considerando o vínculo do usuário, o inventário vigente, os itens esperados e
as ocorrências registradas.

\subsection{Procedimentos executados}

A elaboração da solução e deste artigo foi organizada nas seguintes etapas:

\begin{itemize}
  \item levantamento do problema de conferência patrimonial e delimitação das
  operações que poderiam ser executadas em dispositivo móvel;
  \item análise da aplicação Laravel/Filament existente, dos modelos
  patrimoniais e das tabelas legadas relacionadas a inventário;
  \item definição dos contratos de API e da regra de escopo por usuário
  autenticado;
  \item implementação incremental do backend e do aplicativo Expo/React
  Native;
  \item inspeção dos fluxos implementados para autenticação, consulta, leitura,
  registro de ocorrência e encerramento;
  \item revisão bibliográfica e documental em normas públicas, literatura
  primária e documentação oficial das ferramentas empregadas.
\end{itemize}

\subsection{Requisitos funcionais e não funcionais}

Os requisitos identificados para o recorte de conferência são apresentados na
Tabela \ref{tab:requisitos}.

\begin{table}[H]
\centering
\caption{Requisitos principais da solução}
\label{tab:requisitos}
\small
\begin{tabularx}{\textwidth}{L{1.3cm} Y Y}
\toprule
\textbf{ID} & \textbf{Requisito} & \textbf{Justificativa} \\
\midrule
RF01 & Autenticar o servidor no aplicativo e emitir token de acesso. &
Identificar autoria e proteger endpoints. \\
RF02 & Resolver unidade e setor no backend. &
Impedir que o cliente escolha livremente o escopo patrimonial. \\
RF03 & Listar bens elegíveis do setor e apresentar indicadores. &
Orientar a vistoria e acompanhar progresso. \\
RF04 & Consultar ou ler código patrimonial por digitação ou câmera. &
Atender etiquetas legíveis e situações em que a leitura óptica falhe. \\
RF05 & Registrar bem localizado, não localizado ou divergente. &
Representar os resultados reais da vistoria com evidência textual quando
necessário. \\
RF06 & Finalizar somente atividade sem pendências impeditivas. &
Evitar encerramento incompleto do setor. \\
RNF01 & Proteger comunicação em produção com HTTPS e token. &
Reduzir exposição de credenciais e dados administrativos. \\
RNF02 & Executar gravações relacionadas em transação. &
Manter consistência entre tabelas legadas e itens do inventário. \\
RNF03 & Manter compatibilidade com o banco e o painel já existentes. &
Permitir adoção incremental sem substituir o EGAP desktop. \\
\bottomrule
\end{tabularx}
\end{table}

\subsection{Fontes de evidência e limitações metodológicas}

A descrição técnica foi fundamentada no código-fonte do backend, nas rotas da
API, nos serviços de domínio, nos modelos Eloquent, no aplicativo mobile e nos
arquivos de dependências do projeto \cite{egapRepositorio}. As tecnologias e
conceitos foram confrontados com documentação oficial e fontes normativas ou
primárias.

Nesta versão, não foram realizados testes de usabilidade controlados,
comparação temporal entre inventário manual e mobile ou implantação
experimental acompanhada por métricas. Assim, expressões como redução de
retrabalho indicam finalidade e potencial da solução, e não resultado
quantitativo comprovado.

\section{Desenvolvimento da Solução}

\subsection{Visão arquitetural}

A solução é organizada em três camadas principais:

\begin{itemize}
  \item \textbf{Aplicativo mobile}: interface de campo desenvolvida com Expo,
  React Native, Expo Router e TypeScript; captura ações do servidor e apresenta
  retorno da API;
  \item \textbf{API e regras de aplicação}: aplicação Laravel que autentica,
  resolve escopo, valida operações, aplica regras de inventário e devolve
  respostas JSON;
  \item \textbf{Persistência institucional}: base MySQL/MariaDB já utilizada
  pelo EGAP, contendo bens, inventários, atividades e registros da conferência.
\end{itemize}

\begin{figure}[H]
\centering
\caption{Arquitetura lógica da conferência mobile}
\label{fig:arquitetura}
\fbox{\begin{minipage}{0.92\textwidth}
\centering
\textbf{Servidor em campo}\\
Aplicativo Expo / React Native\\
Login, consulta, câmera, lista, registro e finalização\\[0.25cm]
$\downarrow$ \quad HTTPS + JSON + token Sanctum \quad $\uparrow$\\[0.25cm]
\textbf{API Laravel / Service Layer}\\
\texttt{MobileAuthController} \quad \texttt{ConferenciaBensController}\\
\texttt{UsersConnectionService} \quad \texttt{ConferenciaBensService}\\[0.25cm]
$\downarrow$ \quad Eloquent / transações \quad $\uparrow$\\[0.25cm]
\textbf{Base relacional EGAP}\\
\texttt{mat\_patrimonio} \quad \texttt{mat\_inventario} \quad
\texttt{inv\_atividades} \quad \texttt{mat\_itensinventario}
\end{minipage}}
\end{figure}

Essa arquitetura reduz o acoplamento do cliente ao legado: mudanças de
apresentação mobile não precisam reproduzir regras de inventário, enquanto a
API pode continuar operando sobre as estruturas institucionais existentes.

\subsection{Tecnologias empregadas}

A Tabela \ref{tab:tecnologias} apresenta as tecnologias identificadas nos
arquivos \texttt{composer.json} e \texttt{package.json} da implementação. O
backend utiliza PHP, linguagem documentada pelo PHP Group \cite{phpManual}, e
o painel web utiliza Filament 3.x, cujos painéis reúnem páginas, recursos,
formulários, tabelas e widgets sobre o ecossistema Laravel \cite{filamentDocs}.

\begin{table}[H]
\centering
\caption{Tecnologias empregadas no projeto}
\label{tab:tecnologias}
\small
\begin{tabularx}{\textwidth}{L{2.0cm} L{3.4cm} L{2.0cm} Y}
\toprule
\textbf{Camada} & \textbf{Tecnologia} & \textbf{Versão} & \textbf{Finalidade} \\
\midrule
Backend & PHP & \texttt{\^{}8.2} & Linguagem de execução do servidor. \\
Backend & Laravel & \texttt{\^{}11.0} & Framework web, rotas, validação,
ORM e API. \\
Backend & Filament & \texttt{3.3} & Painel administrativo web do EGAP. \\
Backend & Laravel Sanctum & \texttt{\^{}4.0} & Autenticação por token da API
mobile. \\
Dados & MySQL/MariaDB & --- & Persistência relacional da conexão
\texttt{egap}. \\
Mobile & Expo & \texttt{\~{}54.0.33} & Plataforma e módulos nativos do
aplicativo. \\
Mobile & React Native & \texttt{0.81.5} & Interface móvel multiplataforma. \\
Mobile & TypeScript & \texttt{\~{}5.9.2} & Tipagem da interface e dos
contratos. \\
Mobile & Expo Router & \texttt{\~{}6.0.23} & Navegação baseada em arquivos. \\
Mobile & \texttt{expo-camera} & \texttt{\~{}17.0.10} & Leitura de códigos
patrimoniais. \\
Mobile & \texttt{expo-secure-store} & \texttt{\~{}15.0.8} & Persistência do
token em dispositivos móveis. \\
\bottomrule
\end{tabularx}
\end{table}

\subsection{Modelo de dados da conferência}

A solução reutiliza tabelas patrimoniais existentes. Essa decisão preserva o
histórico e a compatibilidade com o painel administrativo, mas requer que a
API respeite nomenclaturas e estados já consolidados no legado.

\begin{table}[H]
\centering
\caption{Entidades centrais para a conferência}
\label{tab:entidades}
\small
\begin{tabularx}{\textwidth}{L{3.5cm} Y Y}
\toprule
\textbf{Tabela} & \textbf{Conteúdo relevante} & \textbf{Papel no fluxo mobile} \\
\midrule
\texttt{mat\_patrimonio} & Cadastro principal do bem, patrimônio, descrição,
localização, situação e campos de inventário. & Origem dos bens esperados e
alvo da atualização compatível com o legado. \\
\texttt{mat\_inventario} & Ciclo de inventário, número, ano, início, término
e situação. & Determina o inventário considerado atual pela API. \\
\texttt{inv\_atividades} & Atividade do inventário por unidade e setor,
situação e quantidade inventariada. & Controla abertura, edição e conclusão da
conferência setorial. \\
\texttt{mat\_itensinventario} & Registro individual conferido, estado,
observação, responsável e dados copiados do bem. & Armazena localização, não
localização e divergência. \\
\bottomrule
\end{tabularx}
\end{table}

O serviço consulta bens elegíveis por unidade e setor e aceita situações
patrimoniais identificadas no código pelos valores \texttt{1}, \texttt{7} e
\texttt{8}. Também considera estados de transferência já representados no
legado. O item de inventário mantém uma fotografia de informações do bem
durante a conferência, como número patrimonial, descrição, marca, modelo,
estado de conservação e observação, o que favorece posterior auditoria da
ocorrência.

\subsection{Autenticação e determinação do escopo}

O endpoint público \texttt{POST /mobile-api/login} recebe credenciais. O
\texttt{MobileAuthController} tenta autenticar, primeiro, usuários da tabela
local \texttt{users} e, quando necessário, usuários do ambiente EGAP legado.
A senha é verificada com o mecanismo de hash do Laravel. Após encontrar um
vínculo válido, o backend emite um token Sanctum denominado
\texttt{mobile-app}.

O \texttt{UsersConnectionService} relaciona o usuário autenticado a dados
institucionais por CPF e consulta sua lotação mais recente. A partir disso, a
API obtém \texttt{id\_egap}, unidade judiciária e setor. Em todas as rotas de
consulta e conferência, esses dados são reconstruídos no servidor a partir do
token; portanto, a aplicação mobile não pode escolher outro setor apenas
alterando o corpo da requisição.

No cliente, o token é inserido no cabeçalho
\texttt{Authorization: Bearer <token>} de cada requisição autenticada. Em
Android e iOS, o app persiste o token com \texttt{expo-secure-store}; no
logout, a API revoga o token atual e o aplicativo remove a sessão local.

\subsection{Contratos da API mobile}

As rotas de conferência e consulta foram agrupadas sob
\texttt{/mobile-api}. A Tabela \ref{tab:rotas} resume os contratos relevantes
ao recorte deste trabalho.

\begin{longtable}{L{1.45cm} L{4.8cm} L{7.55cm}}
\caption{Endpoints utilizados na conferência patrimonial}
\label{tab:rotas}\\
\toprule
\textbf{Método} & \textbf{Endpoint} & \textbf{Responsabilidade} \\
\midrule
\endfirsthead
\toprule
\textbf{Método} & \textbf{Endpoint} & \textbf{Responsabilidade} \\
\midrule
\endhead
\bottomrule
\endfoot
POST & \texttt{/login} & Autenticar usuário e retornar token e dados do
servidor. \\
GET & \texttt{/me} & Validar sessão autenticada e atualizar identificação do
usuário. \\
POST & \texttt{/logout} & Revogar o token em uso. \\
GET & \texttt{/dashboard} & Obter indicadores patrimoniais e resumo da
conferência atual. \\
GET & \texttt{/bens} & Listar, pesquisar e paginar bens elegíveis do setor. \\
GET & \texttt{/bens/\{codigo\}} & Consultar detalhes de um patrimônio,
inclusive para leitura no painel. \\
GET & \texttt{/conferencia/atual} & Obter inventário, atividade setorial e
métricas. \\
GET & \texttt{/conferencia/bens} & Listar bens esperados acompanhados de
estado de conferência. \\
POST & \texttt{/conferencia/validar-leitura} & Normalizar e validar o código
lido ou digitado antes de qualquer gravação. \\
POST & \texttt{/conferencia/localizar} & Confirmar localização do bem em
transação. \\
POST & \texttt{/conferencia/nao-localizados} & Registrar bens não encontrados,
exigindo justificativa. \\
POST & \texttt{/conferencia/divergencias} & Registrar inconsistência do bem
ou bem físico sem cadastro, exigindo observação. \\
POST & \texttt{/conferencia/finalizar} & Finalizar a atividade quando não
restarem pendências impeditivas. \\
\end{longtable}

\subsection{Interface e fluxo da conferência}

O aplicativo possui rotas para login, painel patrimonial, listagem de bens e
conferência. No painel, o servidor visualiza indicadores de bens e valores,
consulta patrimônios manualmente ou por câmera e mantém histórico local das
consultas recentes. A lista de bens oferece pesquisa e paginação para o escopo
do usuário.

A tela de conferência representa o caso de uso principal. Ao ser aberta, ela
obtém o inventário atual, a atividade e os bens esperados. A interface
apresenta métricas de total, localizados, pendentes, não localizados,
divergentes, em transferência e cadastrados manualmente. Em seguida, o fluxo
operacional ocorre conforme as etapas abaixo:

\begin{enumerate}
  \item o servidor lê a etiqueta pela câmera ou digita o número patrimonial;
  \item o aplicativo envia o código para validação no backend;
  \item a API verifica cadastro, setor, situação, transferência e eventual
  registro prévio;
  \item quando permitido, o servidor confirma localização ou registra uma
  ocorrência acompanhada de justificativa ou observação;
  \item o backend grava o resultado e recalcula o resumo da atividade;
  \item concluída a vistoria, a finalização somente é aceita caso não existam
  pendências impeditivas.
\end{enumerate}

\begin{table}[H]
\centering
\caption{Resultados de leitura tratados pelo backend}
\label{tab:leitura}
\small
\begin{tabularx}{\textwidth}{L{3.6cm} Y}
\toprule
\textbf{Resultado} & \textbf{Interpretação operacional} \\
\midrule
\texttt{localizavel} & Bem do setor, elegível e disponível para confirmação
de localização. \\
\texttt{ja\_conferido} & Bem ou código sem cadastro já possui apontamento no
inventário corrente. \\
\texttt{outro\_setor} & Bem existe, mas não pertence ao escopo setorial do
servidor. \\
\texttt{nao\_cadastrado} & Código físico não corresponde a bem localizado no
cadastro; pode originar divergência documentada. \\
\texttt{situacao\_nao\_conferivel} & Bem possui estado que impede confirmação
ordinária. \\
\texttt{em\_transferencia} & Bem está em fluxo de transferência e não deve ser
localizado automaticamente. \\
\texttt{cadastrado\_manualmente} & Bem exige atenção adicional, embora a API
permita continuidade controlada. \\
\bottomrule
\end{tabularx}
\end{table}

\subsection{Regras de integridade e rastreabilidade}

A confirmação de localização é realizada em transação. O serviço recupera o
bem com bloqueio para atualização, verifica a existência de item prévio, cria
o registro em \texttt{mat\_itensinventario}, atualiza os campos de inventário
em \texttt{mat\_patrimonio} e recalcula a quantidade inventariada da
atividade. Esse fluxo diminui o risco de confirmações concorrentes gerarem
itens duplicados.

Para o estado não localizado, a justificativa é obrigatória tanto na validação
da requisição quanto no serviço. Para divergência, a observação é obrigatória,
podendo informar campos divergentes. O serviço também registra como divergente
um código encontrado fisicamente sem cadastro associado, mantendo o código e a
observação sem inventar uma entidade patrimonial inexistente.

A finalização consulta o resumo calculado no servidor. Havendo bens pendentes,
a API devolve conflito e preserva a atividade em aberto. Após finalização,
novas operações de escrita são bloqueadas. Assim, o aplicativo não determina
sozinho se o inventário terminou: a decisão decorre do estado persistido e das
regras centrais do backend.

\section{Resultados Preliminares e Discussão}

\subsection{Funcionalidades implementadas}

A inspeção da implementação permite identificar os seguintes resultados
funcionais já presentes:

\begin{itemize}
  \item autenticação mobile integrada a usuário local ou vínculo EGAP legado,
  com token Sanctum e validação de lotação;
  \item painel mobile com identificação do escopo, indicadores patrimoniais,
  consulta individual e leitura por câmera;
  \item lista paginada dos bens elegíveis do setor com busca por informações
  patrimoniais;
  \item tela de conferência com métricas, filtros por estado, leitura manual ou
  óptica e ações operacionais;
  \item API dedicada para validação, localização, não localização,
  divergência e finalização;
  \item centralização das regras de inventário em serviço Laravel e persistência
  transacional nas tabelas já utilizadas pelo EGAP.
\end{itemize}

Comparada a um processo em que a conferência física é separada do lançamento
digital, a solução fornece meios técnicos para registrar ocorrências durante a
vistoria e obter retorno imediato da validação. Esse é o principal resultado
preliminar: uma interface de campo integrada ao sistema institucional, sem
transferir para o dispositivo decisões sensíveis sobre escopo e consistência
dos dados.

\subsection{Aderência ao problema e contribuições}

O problema formulado envolve mobilidade sem perda de governança. A solução
atende a esse recorte de três maneiras. Primeiro, utiliza a câmera e uma
interface própria para reduzir a distância entre localização física e registro.
Segundo, não confia ao cliente a determinação da unidade ou setor, pois o
escopo é derivado da autenticação. Terceiro, preserva as tabelas e estados do
sistema existente, permitindo que o inventário mobile se integre à operação
administrativa já disponível no EGAP.

Do ponto de vista de arquitetura, o uso de API e camada de serviço torna a
funcionalidade reutilizável por mais de uma interface e facilita testes
posteriores. Do ponto de vista administrativo, justificativas, observações,
autoria e bloqueio após conclusão são elementos de rastreabilidade compatíveis
com a responsabilidade inerente ao patrimônio público.

\subsection{Limitações observadas}

A solução ainda possui limitações relevantes. Não há evidência, nesta etapa, de
testes com servidores em ambiente operacional ou de medição de tempo, erro ou
satisfação. Portanto, ganhos de eficiência precisam ser validados em estudo
posterior. Também não existe sincronização offline; a conferência depende de
conectividade com a API.

Além disso, a aplicação se integra a um esquema legado, com estados e colunas
preexistentes que aumentam a complexidade de validação. A divergência é
registrada como ocorrência, e não altera automaticamente dados mestres do bem,
decisão adequada para controle, mas que exige tratamento administrativo
posterior. Por fim, o armazenamento seguro de token aplica-se aos aplicativos
nativos; eventual uso web precisa considerar controles próprios de segurança.

\subsection{Trabalhos futuros}

Como continuidade, recomenda-se:

\begin{itemize}
  \item implementar testes automatizados para regras de escopo, transação,
  duplicidade e finalização;
  \item realizar avaliação de usabilidade com servidores que executam
  inventários e comparar indicadores antes e depois da adoção;
  \item estudar modo offline com fila local e sincronização auditável;
  \item incorporar evidências adicionais, como imagens de divergências, sob
  regras institucionais de armazenamento e proteção de dados;
  \item avaliar alertas e painéis gerenciais para acompanhamento da conclusão
  por unidade e setor.
\end{itemize}

\section{Conclusão}

Este artigo apresentou o desenvolvimento de um aplicativo mobile para
conferência patrimonial integrado ao EGAP. A proposta partiu de uma dificuldade
prática: a conferência é realizada onde os bens estão localizados, enquanto o
registro eletrônico pode depender de processos posteriores e manuais. Para
enfrentar essa lacuna, foi implementada uma arquitetura formada por aplicativo
Expo/React Native, API Laravel autenticada por Sanctum e persistência
relacional compatível com as tabelas existentes.

A solução trata funcionalidades específicas do inventário: identificar o
servidor e seu escopo, consultar bens, ler códigos, registrar localização, não
localização ou divergência, atualizar métricas e impedir finalização com
pendências. A análise do código evidencia que regras sensíveis foram
centralizadas no backend e que operações críticas utilizam transações, escolhas
coerentes com rastreabilidade e integridade exigidas no domínio patrimonial
público.

Como resultado preliminar, o Inventário Mobile constitui uma base funcional
para tornar a conferência mais próxima do ambiente real da vistoria e mais
integrada ao EGAP. A comprovação de benefícios operacionais, contudo, depende
de testes, implantação acompanhada e métricas com usuários reais. Desse modo,
o trabalho estabelece o núcleo técnico da solução e define um caminho de
evolução responsável para a digitalização da conferência patrimonial.

\clearpage
\addcontentsline{toc}{section}{Referências}
\begin{thebibliography}{99}

\bibitem{bass2021}
BASS, Len; CLEMENTS, Paul; KAZMAN, Rick.
\textbf{Software Architecture in Practice}. 4. ed. Boston: Addison-Wesley,
2021.

\bibitem{brasil4320}
BRASIL. \textbf{Lei n.º 4.320, de 17 de março de 1964}. Estatui normas gerais
de direito financeiro para elaboração e controle dos orçamentos e balanços.
Disponível em:
\url{https://www.planalto.gov.br/ccivil_03/leis/l4320compilado.htm}.
Acesso em: 25 maio 2026.

\bibitem{cfcnbctsp07}
CONSELHO FEDERAL DE CONTABILIDADE. \textbf{NBC TSP 07: Ativo Imobilizado}.
Brasília: CFC, 2017. Disponível em:
\url{https://portalrestore.cfc.org.br/tecnica/normas-brasileiras-de-contabilidade/nbc-tsp-do-setor-publico/}.
Acesso em: 25 maio 2026.

\bibitem{codd1970}
CODD, Edgar F. A relational model of data for large shared data banks.
\textbf{Communications of the ACM}, v. 13, n. 6, p. 377--387, 1970.
Disponível em:
\url{https://research.ibm.com/publications/a-relational-model-of-data-for-large-shared-data-banks}.
Acesso em: 25 maio 2026.

\bibitem{egapRepositorio}
EGAP. \textbf{Código-fonte da aplicação Laravel e do aplicativo Inventário
Mobile}. Repositório local do projeto, versão analisada em 25 maio 2026.

\bibitem{expoCamera}
EXPO. \textbf{Camera: Expo Documentation}. Disponível em:
\url{https://docs.expo.dev/versions/latest/sdk/camera/}.
Acesso em: 25 maio 2026.

\bibitem{expoRouter}
EXPO. \textbf{Introduction to Expo Router}. Disponível em:
\url{https://docs.expo.dev/router/introduction/}.
Acesso em: 25 maio 2026.

\bibitem{expoSecureStore}
EXPO. \textbf{SecureStore: Expo Documentation}. Disponível em:
\url{https://docs.expo.dev/versions/latest/sdk/securestore/}.
Acesso em: 25 maio 2026.

\bibitem{fielding2000}
FIELDING, Roy Thomas. \textbf{Architectural Styles and the Design of
Network-based Software Architectures}. 2000. Tese (Doutorado) -- University of
California, Irvine, 2000. Disponível em:
\url{https://ics.uci.edu/~fielding/pubs/dissertation/fielding_dissertation.pdf}.
Acesso em: 25 maio 2026.

\bibitem{filamentDocs}
FILAMENT. \textbf{Getting Started: Filament 3.x Panels Documentation}.
Disponível em:
\url{https://filamentphp.com/docs/3.x/panels/getting-started}.
Acesso em: 25 maio 2026.

\bibitem{fowlerServiceLayer}
FOWLER, Martin. \textbf{Service Layer}. In: Patterns of Enterprise Application
Architecture. Disponível em:
\url{https://martinfowler.com/eaaCatalog/serviceLayer.html}.
Acesso em: 25 maio 2026.

\bibitem{laraveldatabase}
LARAVEL. \textbf{Database: Getting Started -- Laravel 11.x Documentation}.
Disponível em: \url{https://laravel.com/docs/11.x/database}.
Acesso em: 25 maio 2026.

\bibitem{laraveleloquent}
LARAVEL. \textbf{Eloquent ORM -- Laravel 11.x Documentation}. Disponível em:
\url{https://laravel.com/docs/11.x/eloquent}.
Acesso em: 25 maio 2026.

\bibitem{laravelrouting}
LARAVEL. \textbf{Routing -- Laravel 11.x Documentation}. Disponível em:
\url{https://laravel.com/docs/11.x/routing}.
Acesso em: 25 maio 2026.

\bibitem{laravelsanctum}
LARAVEL. \textbf{Laravel Sanctum -- Laravel 11.x Documentation}. Disponível em:
\url{https://laravel.com/docs/11.x/sanctum}.
Acesso em: 25 maio 2026.

\bibitem{mysqlForeignKeys}
MYSQL. \textbf{MySQL 8.4 Reference Manual: FOREIGN KEY Constraints}.
Disponível em:
\url{https://dev.mysql.com/doc/refman/8.4/en/create-table-foreign-keys.html}.
Acesso em: 25 maio 2026.

\bibitem{mysqlTransactions}
MYSQL. \textbf{MySQL 8.4 Reference Manual: InnoDB Transaction Model}.
Disponível em:
\url{https://dev.mysql.com/doc/refman/8.4/en/innodb-transaction-model.html}.
Acesso em: 25 maio 2026.

\bibitem{phpManual}
PHP GROUP. \textbf{PHP Manual: Documentation}. Disponível em:
\url{https://www.php.net/docs.php}.
Acesso em: 25 maio 2026.

\bibitem{reactnativeDocs}
REACT NATIVE. \textbf{Introduction: React Native Documentation}. Disponível
em: \url{https://reactnative.dev/docs/getting-started}.
Acesso em: 25 maio 2026.

\bibitem{rfc8446}
RESCORLA, Eric. \textbf{The Transport Layer Security (TLS) Protocol Version
1.3}. RFC 8446. Internet Engineering Task Force, 2018. Disponível em:
\url{https://www.rfc-editor.org/rfc/rfc8446}.
Acesso em: 25 maio 2026.

\bibitem{rfc9110}
FIELDING, Roy; NOTTINGHAM, Mark; RESCHKE, Julian.
\textbf{HTTP Semantics}. RFC 9110. Internet Engineering Task Force, 2022.
Disponível em: \url{https://www.rfc-editor.org/rfc/rfc9110}.
Acesso em: 25 maio 2026.

\end{thebibliography}

\end{document}
